%% =====================================================================
%%  T-19-03  The Unanswerable Question
%%  -------------------------------------------------------------------
%%  One idea per file. Core is mandatory. Slots in canonical order.
%%  Max 700 words. No layout commands. No filler. New source material
%%  for this entry goes in sources/B1/T-19-03.tex -- never by
%%  rewriting this file (docs/RULES.md R0.1).
%%  =====================================================================
%% @kind: topic
%% @id: T-19-03
%% @title: The Unanswerable Question
%% @subtitle: Pressure without a raised voice --- the audience does the coercing
%% @chapter: c19
%% @order: 30
%% @status: stable
%% @sources: B1
%% @updated: 2026-09-19

\topic{T-19-03}{The Unanswerable Question}{Pressure without a raised voice --- the audience does the coercing}

\begin{core}
The gentlest member of the pressure family: ask a person, in front of onlookers, a question that cannot be answered with no. The question itself is harmless; the audience converts it into obligation, because refusing in public costs a self-image. B1's strolling minstrel --- \emph{what song do you like?} --- had the author over a barrel before the second verse. This entry documents the device so it can be named and refused.
\end{core}
\begin{mechanism}
The mechanism is audience-as-lever. A private request can be declined at zero cost; the same request, embedded in a question whose grammar excludes refusal and staged before witnesses, attaches a price to declining --- the public label (cheapskate, killjoy, ungenerous, difficult). B1 classes this as one of the working pressures: the target complies not because the want is theirs but because the scene has been arranged so that the alternative costs more than the compliance. The device's cousins in the same chapter: the assumed yes (acting on indecision as if it were agreement) and the silence after the ask (the most powerful pressure short of physical threat, B1 says) --- all three work by making the target's exit more expensive than their consent. The unanswerable question is the social variant, and the most defensible to read: once named, it is almost comically visible.
\end{mechanism}
\begin{conditions}
Strongest in public settings with mixed audiences --- restaurants, meetings, ceremonies, family tables --- and against the image-conscious. It weakens against the shameless (no audience cost, no mechanism), in private (the lever disappears), and among people who have named the device --- the named trick dies on the spot, since the audience's laugh converts the pressure onto the asker.
\end{conditions}
\begin{application}
  \item Recognition first. The signatures: the question offers no grammatical room for no; it is asked where others are watching; the asker's stake appears only after you answer.
  \item Answer the question you wish had been asked: \emph{I'd rather not, thanks} --- delivered pleasantly, it breaks the grammar and the lever together.
  \item Name it lightly if pressed: \emph{that's a question with no wrong answer except honesty}. Public naming flips the embarrassment to its author.
  \item If you are the onlooker, rescue the target --- change the subject, answer for them, applaud their refusal. The audience is the weapon; disarm it.
  \item In your own asking, audit honestly: would this request survive being asked privately? If not, the audience is the point --- and that is this entry, being run on someone.
\end{application}
\begin{feedback}
  \item Working (as defence): the asker rephrases into a private, answerable request. The device has been retired.
  \item Working: the room's sympathy moves to whoever names the mechanism. Naming is the counter-lever.
  \item Failing: you refuse rudely instead of pleasantly --- the label snaps on, and the device has won on a technicality.
  \item Failing: you comply \emph{and} resent it. The worst ledger: full price paid, lesson unbanked.
\end{feedback}
\begin{failure}
The device has a narrow operating window and a steep reversal: audiences enjoy seeing a lever exposed more than they enjoy watching it worked, so the operator risks the room's goodwill every time. Chronic users also train their circles --- people start answering the unanswerable question with visible irony, and the operator's invitations thin out. And the habit generalises: a person who arranges rooms to prevent refusal stops being able to hear refusal at all, which bankrupts every honest channel they have.
\end{failure}
\begin{countermeasures}
  \item Teach the grammar check: before answering any public request, ask whether \emph{no} is a grammatical answer to what was actually asked. If not, the question is the trap.
  \item Practise the pleasant refusal out loud; it is a skill with a rehearsal cost, and the unprepared always pay retail (T-24-05).
  \item In your own meetings, watch for questions designed for the observers. Answer the private version, publicly, and see who squirms.
  \item Notice who seats you beside whom at tables. Stage-managed audiences are this entry, pre-built (T-17-05).
\end{countermeasures}

\sources{B1}
\seesources
\seealso{T-24-05, T-20-01, T-14-03}
